\chapter{{\fontsize{14pt}{21pt}\selectfont\MakeUppercase{АНАЛИЗ ПРЕДМЕТНОЙ ОБЛАСТИ И ОБОСНОВАНИЕ ВЫБОРА ТЕХНОЛОГИЙ}}}
\label{cha:analysis}

\section{Исследование инфраструктуры нагрузочного тестирования и роли имитационных сервисов}

В современной практике обеспечения качества сложных программных комплексов нагрузочное тестирование занимает критически важное место. Оно позволяет гарантировать стабильность и управляемость информационной системы при высоких эксплуатационных нагрузках. Одной из ключевых проблем при организации таких испытаний является необходимость взаимодействия объекта тестирования с множеством внешних систем и сервисов.

Использование реальных внешних интеграций в среде тестирования часто оказывается невозможным или нецелесообразным по ряду причин:

\begin{compactlist}
\item \textbf{Ограниченность ресурсов}: внешние системы могут иметь жесткие лимиты на количество запросов или требовать значительных финансовых затрат на эксплуатацию.
\item \textbf{Изоляция среды}: для получения достоверных результатов необходимо исключить влияние задержек и сбоев сторонних систем на показатели исследуемого объекта.
\item \textbf{Риски изменения данных}: проведение тестов на реальных интеграциях может привести к нежелательному изменению или повреждению информации в продуктовых базах данных.
\end{compactlist}

Для решения этих проблем применяются \textbf{имитационные сервисы («заглушки»)} -- специализированное программное обеспечение, которое воспроизводит поведение реальных компонентов, отвечая на запросы тестируемой системы заранее определенным образом. В контексте данной работы в качестве имитационных сервисов рассматриваются Java-приложения в форматах \textbf{.jar} и \textbf{.war}.

Роль имитационных сервисов в инфраструктуре нагрузочного тестирования заключается в следующем:

\begin{compactlist}
\item \textbf{Обеспечение автономности стенда:} заглушки позволяют полностью изолировать тестируемую систему от внешнего окружения, создавая предсказуемую среду.
\item \textbf{Эмуляция различных сценариев:} имитационные сервисы могут воспроизводить не только корректные ответы, но и ошибки (HTTP 5xx, тайм-ауты), что необходимо для проверки механизмов отказоустойчивости.
\item \textbf{Контроль интенсивности трафика:} использование механизмов ограничения запросов (Rate Limiting) позволяет моделировать поведение внешних систем при достижении ими предельных порогов производительности.
\end{compactlist}

Однако при увеличении масштаба информационной системы количество необходимых заглушек возрастает, что порождает проблему управления распределенной инфраструктурой. Процессы развертывания, контроля состояния (активна/остановлена) и маршрутизации трафика к десяткам имитационных сервисов на различных серверах в ручном режиме существенно замедляют подготовку к испытаниям.

Таким образом, для эффективного функционирования стендов нагрузочного тестирования необходима специализированная система. Она должна централизованно управлять жизненным циклом заглушек, обеспечивать их мониторинг и динамическое проксирование запросов, что особенно актуально в условиях перехода на отечественные операционные системы, такие как \textbf{ОС Astra Linux}.